\section{Theorem Map for DQDock}\label{sec:theorem-map}

For writing purposes, the most useful way to view the formal material is as a map from theorem families to engine responsibilities. Table~\ref{tab:theorem-map} gives the compact version. The point of the table is not merely bookkeeping. It shows where the article's main claims live and how the proof families support the runtime architecture.

\begin{table}[t]
\centering
\small
\begin{tabular}{@{}p{0.12\linewidth}p{0.38\linewidth}p{0.40\linewidth}@{}}
\toprule
\textbf{Handles} & \textbf{Formal role} & \textbf{DQDock consequence} \\
\midrule
\texttt{LS1--4} & explicit lattice-tail control for LJ-type interactions & justifies the error model behind cutoff reasoning \\
\texttt{LJ1--6} & exact-vs-cutoff or softened LJ approximation statements & supports certified or conditionally certified LJ scoring paths \\
\texttt{CB1--4} & Coulomb / real-space Ewald approximation packaging & supports assumption-dependent electrostatics claims \\
\texttt{BP1--3} & gap-plus-tail bridge to bounded-potential locality & turns physical decay control into decision-safe cutoffs \\
\texttt{MD1--7} & outside-cutoff irrelevance and molecular structural-rank bounds & explains pocket-local tractability and thermodynamic interpretation \\
\texttt{APX1--3} & abstract uniform approximation and optimizer invariance & explains when coarse screening preserves exact winners \\
\texttt{SD1--9} & sampled-docking transport of exact/coarse and locality guarantees & supports finite candidate-set docking and restricted support claims \\
\texttt{TK1--6} & pairwise order, ambiguity-band, and top-k preservation & supports ranking-safe filtering and survivor-set reasoning \\
\texttt{CP1} & sound pruning certificate & formal basis for certified elimination of hopeless poses \\
\texttt{GD1--3} & grid/discretization transport from continuous control to finite approximation & supports grid-based or discretized variants of the engine \\
\bottomrule
\end{tabular}
\caption{Handle families added for the docking/computational-chemistry layer of Paper~4 and their intended role in DQDock.}
\label{tab:theorem-map}
\end{table}

Three interpretive lessons follow from this map. First, the formal core is strongest on approximation, locality, and ranking preservation. Second, the map already supports a substantial theorem-backed docking story across the key reductions used by the engine. Third, the code and proof layers line up best when the paper speaks in terms of certified reductions: exact scorer to coarse scorer, ambient coordinates to pocket-local coordinates, large candidate set to certified survivor set.

These reductions suggest a concise paper-facing theorem narrative:
\begin{enumerate}
\item exact physical interaction tails admit explicit approximation control;
\item approximation control plus positive decision gap yields exact winner preservation;
\item winner preservation plus cutoff locality yields pocket-local structural rank;
\item pocket-local structural rank and sampled-docking transport justify restricted candidate families;
\item ranking-preservation and pruning theorems turn those restrictions into certified screening logic.
\end{enumerate}
This five-step progression is, in effect, the formal skeleton of DQDock.

\paragraph{What can be claimed directly.}
The paper can directly claim that DQDock is organized around theorem-backed sufficient conditions for safe cutoff locality, winner preservation under bounded approximation, and certified top-k / pruning behavior. These are concrete support statements for the implemented system. In particular, the structural-rank theorems \leanmeta{\LHrng{MD}{1}{7}.} and the ranking/pruning theorems \leanmeta{\LHrng{TK}{1}{6}, \LH{CP1}.} support a clear formal statement of why the engine is allowed to ignore, coarsen, or discard certain information.

\paragraph{Assumption-dependent layers.}
The theorem map contains packaging theorems as well as direct invariance theorems. The Coulomb and Ewald layer, for example, is theorem-backed once the relevant physical assumptions are supplied. Likewise, the sampled-docking transport results require support and compatibility hypotheses. The article can therefore state the formal architecture positively: the most important approximation and pruning moves in DQDock are controlled by explicit theorem families, while the remaining physics and engineering assumptions are surfaced clearly alongside them.
